% Supplementary Table: 2024 TRIPOD+AI Reporting Audit (Collins et al., BMJ 2024;385:e078378)
\begin{table}[h!]
\centering
\caption{TRIPOD+AI reporting checklist mapping for the retrospective clinical prediction and feature-acquisition evaluations (PhysioNet 2019 ICU Sepsis and Yale New Haven Health System ED Triage; \citealp{collins2024tripodai}).}
\label{tab:tripod_ai}
\scriptsize
\renewcommand{\arraystretch}{1.12}
\begin{tabular}{>{\raggedright\arraybackslash}p{0.12\textwidth}>{\raggedright\arraybackslash}p{0.24\textwidth}>{\raggedright\arraybackslash}p{0.56\textwidth}}
\toprule
\textbf{Section / Item} & \textbf{TRIPOD+AI Requirement} & \textbf{Manuscript Location and Compliance Summary} \\
\midrule
Title (1) & Identify study developing/evaluating a multivariable prediction model using ML & Title specifies threshold-aware feature acquisition and residual learning for clinical prediction; retrospective evaluation explicitly framed in Abstract and Section~\ref{sec:study_design}. \\
Abstract (2) & Structured summary of objectives, data, methods, results, limitations & Narrative $<200$-word Abstract states problem, unified Beta--KL method, synthetic and external hospital cohorts, repaired Common-Panel Bagged VOI results, and retrospective/cost limitations. \\
Background (3a--3b) & Clinical/methodological context, intended role, target threshold & Sections~\ref{sec:intro}--\ref{sec:methods} formulate the Decision-Information Dilution Hypothesis around prespecified clinical threshold $\tau^* = C_{\text{FP}}/(C_{\text{FP}}+C_{\text{FN}})$. \\
Data Sources (4a--4b) & Study design, data sources, dates, eligibility, open-access provenance & Section~\ref{sec:study_design} documents PhysioNet/CinC 2019 ICU cohorts (BIDMC $N=19{,}873$, Emory $N=18{,}359$) and Yale ED cohorts ($N=560{,}486$ across Campuses A, B, and C). \\
Participants (5a--5c) & Setting, inclusion/exclusion criteria, landmark alignment & Causal observation landmarks ($4/6/8/12\text{h}$) exclude pre-landmark \texttt{SepsisLabel}$=1$; measured-bloodwork (Track B1b) vs.\ all-comers (Track A) selection explicitly audited (Appendix~\ref{app:clinical_gov}). \\
Outcome (6a--6c) & Outcome definition, timing/horizon, incorporation-bias assessment & Section~\ref{sec:res_sepsis} and Appendix~\ref{app:horizons_shift} separate ICU observation window from fixed future prediction horizons ($2\text{h}, 6\text{h}, 12\text{h}$, remaining stay) and audit SOFA-constituent incorporation bias. \\
Predictors (7a--7c) & Baseline vs.\ Stage-2 channels, point-in-time availability, cost ledger & Section~\ref{sec:methods} distinguishes bedside Stage-1 predictors from Stage-2 laboratory panels (single-draw \texttt{\_last} vs.\ $6\text{h}$ retrieval summaries) and prices procedure + CPT assay costs via \texttt{PanelCostSchedule}. \\
Sample Size (8) & Cohort sizes, event counts, finite-sample power constraints & Exact $N$ and event counts reported in every table ($349$ external events in Emory Track B1b; $616$ events in Yale Campus C ESI 4--5); split-pilot power loss quantified in Appendix~\ref{app:proofs}. \\
Missing Data (9) & Missingness rates by site/outcome, handling, ordering confounding & Appendix Table~\ref{tab:missingness_avail} reports per-channel availability across sites and outcomes ($Y=0$ vs.\ $Y=1$), complete-panel subset results, and \texttt{AgnielDeconfounder} / \citet{vonkleist2025evaluation} observability assumptions. \\
Analytical Methods (10a--10f) & Learner architectures, hyperparameter tuning, OOF stacking, VOI routing & Section~\ref{sec:methods} and Appendix~\ref{app:proofs} specify LightGBM capacities, 3-fold OOF stacking, empirical martingale VOI offset, and disjoint $\mathcal{D}_{\text{train}}/\mathcal{D}_{\text{cal}}/\mathcal{D}_{\text{test}}$ splits. \\
Fairness / Subgroups (11) & Treatment of sociodemographic attributes and subgroup evaluation & Nine demographic/payer/site variables stripped from fairness-blinded Yale models; $10$ external subgroups audited for Net Benefit (vs.\ Triage and Global-CMI), query rate, sensitivity, false alerts, and \texttt{CITL} (Appendix Table~\ref{tab:fairness}). \\
Class Imbalance \& Shift (12) & Handling of low prevalence and cross-site calibration shift & Unweighted canonical Bernoulli likelihood preserved; zero-shot, unlabeled \citet{saerens2002adjusting} EM, chronological split historical prior, and oracle prevalence adjustments compared (Table~\ref{tab:horizons_shift}). \\
Evaluation Metrics (13a--13c) & Discrimination, calibration, clinical utility, uncertainty intervals & Reports $\operatorname{NB}_{1000}$, $\operatorname{sNB}$, TP/FP/query exchange rates, AUROC, AUPRC, \texttt{CITL}, logistic slope, LOESS \texttt{ICI}, and $B=10{,}000$ paired bootstrap $95\%$ CIs. \\
Development vs.\ Validation (14--17) & Separation of development, sensitivity, and external transport & Provenance of \texttt{v1.0} (\texttt{8c8c4c8}) through \texttt{v2.3.0} disclosed; untested Yale Campus C ($N=71{,}706$ all-comers; $31{,}430$ ESI 4--5) evaluated zero-shot; encounter-level ID limitation of \citet{hong2018predicting} stated. \\
Model Updating / Governance (18--22) & Recalibration and conformal alerting governance & Repeated-subsample ($R=200$) Campus B \texttt{CfBH} burn-in curve across $n_{\text{burn}}\in\{250,\dots,4{,}000\}$ quantifies empirical FDR, recall, and alert volume (Appendix Table~\ref{tab:cfbh_burnin}). \\
Open Science (23--27) & Code, environment lockfile, seeds, protocol registration, funding/COI & MIT-licensed code, \texttt{requirements-lock.txt}, \texttt{SEED\_MANIFEST.md}, checkpoint JSONs, and deterministic table verifier archived at \texttt{github.com/deesatzed/collision-zone-ed} (Release \texttt{v2.3.0}); no external funding or COI; AI workflow disclosed. \\
\bottomrule
\end{tabular}
\end{table}
